\documentclass[11pt]{article}
\usepackage[a4paper,margin=2.6cm]{geometry}
\usepackage[T1]{fontenc}
\usepackage{parskip}
\usepackage[hidelinks]{hyperref}
\title{}
\begin{document}
\begin{center}\large\bfseries Cover Letter --- Cybersecurity (SpringerOpen)\end{center}
\vspace{1em}
Dear Editors of \textit{Cybersecurity},

We are pleased to submit our manuscript entitled ``Evidence-Preserving Graph Denoising for Noisy-Label Intrusion Detection: A Leakage-Audited Study of When Graph Structure Helps'' for consideration as a Research article in \textit{Cybersecurity}.

Intrusion-detection classifiers are routinely trained on delayed, rule-generated, or inconsistent labels, and the standard remedy---dataset purification by deleting suspicious samples---risks discarding exactly the rare, boundary attack evidence that security analysts need most. Our manuscript reframes noisy-label intrusion detection as \emph{evidence-preserving data curation} and reports a leakage-audited, fully reproducible study on three real-world datasets (CICIDS-2017, CESNET-TLS-Year22, UNSW-NB15). We make three contributions: (i) a multi-view graph consistency diagnostic (Graph-CDM) that yields a modest but statistically significant aggregate improvement over the CoLD baseline under a leakage-controlled protocol; (ii) evidence-preserving soft weighting that matches hard deletion on aggregate Macro-F1 while significantly improving rare/tail attack-class Macro-F1 under high asymmetric label noise, benchmarked against external anchors (Confident Learning, Co-Teaching) under an identical protocol with honest reporting of where our approach does not lead; and (iii) a leakage audit showing that naive graph denoising gains can be dramatically inflated by duplicate records and clean-label signal, together with the full reproducibility package.

We believe this work fits the scope of \textit{Cybersecurity} closely, given the journal's focus on data-driven security, intrusion detection, and applied machine learning for security operations.

We confirm that this manuscript is original, has not been published before, and is not under consideration by any other journal. All authors have read and approved the manuscript and agree with its submission to \textit{Cybersecurity}. The authors declare no competing interests, and no specific funding was received for this work. In accordance with the journal's double-anonymous review policy, the main manuscript file contains no identifying information; author details are provided in the separate title page file.

A declaration of generative AI use in the writing process is included in both the manuscript and the title page, in line with Springer Nature policy.

\vspace{1.5em}
Sincerely,

Tianqi Guo (on behalf of all authors)\
University of Science and Technology of China\
guotianqi@mail.ustc.edu.cn
\end{document}
